\documentclass[11pt]{article}
\usepackage[margin=1in]{geometry}
\usepackage[strings]{underscore}
\usepackage{amsmath}
\usepackage{hyperref}
\usepackage{enumitem}
\usepackage{titlesec}
\usepackage{parskip}
\usepackage{xcolor}
\usepackage{microtype}
\sloppy

\hypersetup{colorlinks=true, linkcolor=blue!50!black, urlcolor=blue!50!black, citecolor=blue!50!black}
\titleformat{\section}{\normalfont\Large\bfseries}{\thesection}{1em}{}
\titleformat{\subsubsection}[block]{\normalfont\bfseries\ttfamily\footnotesize}{}{0pt}{}

\newcommand{\role}[1]{\par\noindent\textit{#1}\par\vspace{0.5em}}
\newcommand{\func}[1]{\subsubsection*{\texttt{#1}}}
\newcommand{\cliarg}[2]{\item \texttt{#1} --- #2}

\title{Bruise Segmentation Pipeline\\\large Function-by-Function Reference}
\author{BRUISE\_SEGMENTATION\_PROJECT}
\date{Generated for the 5-model scope: SegFormer B2 teacher, SegFormer B0 direct,\\
SegFormer B0 distilled, YOLO26n-sem direct, YOLO26n-sem distilled}

\begin{document}
\maketitle

\section*{Scope of this document}
This reference covers every file remaining in \texttt{pipeline/} and \texttt{scripts/}
after the ALS / fairness-distillation / ITA / skin-tone files were moved to
\texttt{out\_of\_scope/} (they are not documented here, since current project focus
is only the 5 models listed above). Part~I covers the shared \texttt{pipeline/}
modules; Part~II covers the numbered pipeline scripts and standalone utility scripts
in \texttt{scripts/}, in the order they are actually run.

\tableofcontents
\newpage

\part{Shared Pipeline Modules (\texttt{pipeline/})}

\section{pipeline/data.py}
\role{Defines the dataset class and preprocessing transform used by every
training, calibration, threshold-search, and evaluation step in the project ---
for both SegFormer and YOLO. Centralizing this in one file guarantees that
whatever operating point (threshold, temperature) gets calibrated on the
validation set is always applied to the same input distribution at test time.}

\func{\_find\_col(df: pd.DataFrame, candidates: list[str]) -> str | None}
Searches a manifest DataFrame's columns for the first name in \texttt{candidates}
that exists, trying an exact match first and then a case-insensitive match.
Exists because different manifest CSVs in this project use slightly different
column names for the same concept (e.g. \texttt{mask\_path} vs.
\texttt{majority\_mask\_path}), and this lets one normalization function handle
all of them without hardcoding a single schema.

\func{normalize\_manifest(df: pd.DataFrame) -> pd.DataFrame}
Standardizes any manifest DataFrame to guarantee \texttt{image\_path},
\texttt{mask\_path}, \texttt{stem}, and \texttt{subject} columns exist, deriving
\texttt{stem} from the image filename and \texttt{subject} from the stem (splitting
on underscore/hyphen) if not already present. Raises \texttt{RuntimeError} with the
full list of available columns if no image or mask column can be found ---
failing loudly immediately rather than crashing confusingly deep inside a
training loop. Every downstream script calls this before touching a manifest,
so schema differences between manifests are handled in exactly one place.

\func{load\_train\_val\_split(csv\_path: str) -> tuple[pd.DataFrame, pd.DataFrame]}
Loads and normalizes a split CSV, then splits it into train/val DataFrames based
on its \texttt{split} column, raising \texttt{RuntimeError} if that column is
missing (telling the caller to run \texttt{00\_build\_split.py} first).

\func{load\_fixed\_test(csv\_path: str) -> pd.DataFrame}
Loads and normalizes the fixed held-out test manifest. A thin wrapper so every
script that reads the test set goes through the same normalization path as
train/val.

\func{get\_augmentation(training: bool, img\_h: int, img\_w: int) -> A.Compose}
Builds the Albumentations transform pipeline: resize to \texttt{img\_h} x
\texttt{img\_w}, ImageNet mean/std normalization, convert to tensor. When
\texttt{training=True}, adds horizontal/vertical flips, brightness/contrast
jitter, hue/saturation/value jitter, and Gaussian noise; when \texttt{False}
(validation, test, and every benchmark/inference use in this project), it is the
deterministic resize+normalize+to-tensor pipeline only. This is the single
function reused by \texttt{pipeline/benchmark\_640.py} for both SegFormer and
YOLO inference, guaranteeing the benchmark sees the same input distribution the
model was calibrated on.

\func{class BruiseDataset(Dataset)}
The shared PyTorch \texttt{Dataset} implementation used by every SegFormer
training/eval script and by the YOLO threshold/temperature search.
\texttt{\_\_init\_\_(df, img\_h, img\_w, training=False)} stores the manifest
DataFrame and builds its transform via \texttt{get\_augmentation}.
\texttt{\_\_len\_\_()} returns the row count. \texttt{\_\_getitem\_\_(idx)} reads
the image with OpenCV, converts BGR to RGB, reads the mask as grayscale and
binarizes it as \texttt{mask > 0} (not \texttt{mask == 255}) --- robust to masks
whose foreground pixel value is not exactly 255 (e.g. anti-aliased or
class-index masks) --- applies the transform to both image and mask jointly (so
any spatial augmentation stays aligned between them), and returns
\texttt{(image\_tensor, mask\_tensor, stem, image\_path, mask\_path)}. Returning
the stem and paths alongside the tensors lets every evaluation script attribute
a metric row back to a specific source image without a second lookup.

\newpage
\section{pipeline/io\_utils.py}
\role{Foundation utilities shared by every script: YAML/JSON I/O, directory
creation, structured logging, and pre-flight config validation. Centralizing
these means every script fails with an identical, clear error message for a
missing config or bad path, rather than each script inventing its own
(inconsistent) error handling.}

\func{load\_yaml(path: str | Path) -> dict}
Loads a YAML file, raising \texttt{FileNotFoundError} with an actionable message
(run from the project root, or pass an absolute path) if it does not exist ---
rather than the bare, less-clear error \texttt{open()} itself would raise.
Coerces an empty file to \texttt{\{\}} so callers can unconditionally use
\texttt{cfg.get(...)} without an \texttt{AttributeError}.

\func{ensure\_dir(path: str | Path) -> Path}
Creates a directory and all missing parents (\texttt{parents=True,
exist\_ok=True}), returning the \texttt{Path} so callers can chain it, e.g.
\texttt{csv\_path = ensure\_dir(run\_dir / "eval") / "results.csv"}.

\func{save\_json(path: str | Path, payload: dict) -> None}
Serializes a dict to a pretty-printed (\texttt{indent=2}) JSON file, creating
parent directories first. Used for \texttt{run\_config.json} and
\texttt{temperature.json} so downstream scripts can read back exact
hyperparameters without re-parsing CLI arguments.

\func{setup\_logging(run\_dir: Path | str | None = None, level: int = logging.INFO) -> logging.Logger}
Configures and returns the shared \texttt{"pipeline"} logger with a consistent
timestamped format, writing to stdout always and additionally to
\texttt{run\_dir/run.log} if a \texttt{run\_dir} is given (append mode, so a
resumed run keeps its history). Clears any existing handlers first so calling
this twice in one process never duplicates log lines.

\func{validate\_paths(paths: dict) -> None}
Pre-flight check for \texttt{configs/paths.yaml}: asserts every required key is
present (collecting \emph{all} missing keys before raising, not just the first,
so a user fixes every typo in one pass), then checks that every required
\emph{input} path (manifests, pretrained weights, YOLO weights) actually exists
on disk, raising \texttt{RuntimeError} listing every missing path at once.
Output-only paths (like \texttt{project\_root}) are not checked here since
scripts create those themselves.

\func{validate\_cfg(cfg: dict) -> None}
Pre-flight check for \texttt{configs/common\_train.yaml}: asserts every required
key is present, then sanity-checks numeric ranges (e.g. \texttt{epochs} in
$[1, 10000]$, \texttt{patience} $\le$ \texttt{epochs}, learning rates in $(0,1)$,
a non-empty \texttt{thresholds} list) so a unit-error typo (like a learning rate
of 6 instead of 0.00006) is caught before any GPU time is spent, not discovered
hours into a training run.

\newpage
\section{pipeline/models.py}
\role{Defines the SegFormer model wrapper and the checkpoint/threshold loading
functions used by every SegFormer training and evaluation script (including
both benchmark scripts).}

\func{build\_segformer(pretrained: str, num\_labels: int = 1) -> nn.Module}
Loads a HuggingFace \texttt{SegformerForSemanticSegmentation} from a local
pretrained directory, with \texttt{num\_labels=1} (binary bruise segmentation)
and \texttt{ignore\_mismatched\_sizes=True} (since the pretrained checkpoint's
original 1000-class ImageNet decode head does not match a 1-class head --- that
head is discarded and randomly reinitialized, while the backbone/encoder weights
load normally).

\func{class SegformerWrapper(nn.Module)}
Wraps the HuggingFace model, exposing \texttt{.backbone} and
\texttt{.decode\_head} as properties so \texttt{pipeline/scheduler.py} can give
them different learning rates. \texttt{forward(x)} runs the model and, if the
returned logits are not already at the input's spatial resolution, upsamples
them (bilinear) to match --- this is why \texttt{SegFormer640Adapter} in
\texttt{benchmark\_640.py} never needs a separate resize step.
\texttt{gradient\_checkpointing\_enable()} turns on gradient checkpointing on the
encoder if available, silently doing nothing if the underlying HuggingFace
version does not expose that attribute.

\func{count\_params(model: nn.Module) -> tuple[int, int]}
Returns \texttt{(total\_parameters, trainable\_parameters)} by summing
\texttt{.numel()} over all parameters (the second sum filtered to
\texttt{requires\_grad}).

\func{load\_segformer\_threshold(run\_dir: Path) -> float}
Reads a SegFormer run's \texttt{threshold\_search.csv} (written by the training
loop's validation-set sweep), sorts by \texttt{mean\_dice} descending, and
returns the top row's threshold. This is always read from the \textbf{validation}
set, never recomputed on test --- fitting a threshold on the same data used to
report accuracy would leak information and inflate the reported score.

\func{\_pick\_matching\_checkpoint(run\_dir: Path, model: nn.Module) -> tuple[Path, dict]}
Returns whichever of \texttt{best\_model.pt} or
\texttt{best\_model.pre\_migration\_backup.pt} actually has parameter names
matching the currently-installed \texttt{transformers} version. This exists
because \texttt{26\_migrate\_legacy\_segformer\_checkpoints.py} rewrote some
checkpoints' internal parameter names to match the \texttt{transformers} version
ORC had at one point, keeping the pre-migration original as a safety backup ---
but a different machine's \texttt{transformers} install may match one file or
the other, so both are loaded and compared by key set (not by tensor values,
which is cheap since both are already read into memory) before either is
actually used. Raises \texttt{RuntimeError} naming both candidates if neither
matches, pointing at script 26's docstring for how to add a new remapping rule.

\func{load\_segformer\_model(model\_name: str, pretrained\_key: str, paths: dict, device: torch.device) -> tuple[nn.Module, float, Path]}
The single function every SegFormer-consuming script (training scripts'
\texttt{load\_teacher}, script 29, and \texttt{benchmark\_640.py}) uses to load a
trained run: resolves the threshold via \texttt{load\_segformer\_threshold},
builds a fresh model via \texttt{build\_segformer}/\texttt{SegformerWrapper}, and
loads whichever checkpoint file \texttt{\_pick\_matching\_checkpoint} determines
is compatible. Returns \texttt{(model, threshold, checkpoint\_path)} so callers
can both run inference and record exactly which checkpoint produced a given
result.

\newpage
\section{pipeline/batch\_finder.py}
\role{Automatically probes the largest training batch size that fits in
available GPU memory, so a single hardcoded batch size does not under-utilize
the GPU for a small model or crash with out-of-memory on a large one (e.g. the
teacher, which additionally must carry a frozen teacher model through every
forward pass during distillation).}

\func{find\_optimal\_micro\_batch(model: nn.Module, img\_h: int, img\_w: int, device: torch.device, effective\_batch: int = 8, target\_hi: float = 0.75, amp: bool = True, max\_probe: int = 32, teacher\_fn=None) -> tuple[int, int, float]}
Probes increasing batch sizes (doubling each time: 1, 2, 4, 8, \ldots) with a
real forward + backward + optimizer step on random dummy data at the model's
actual training resolution, stopping when either an out-of-memory error occurs
or the peak VRAM fraction exceeds \texttt{target\_hi} (default 75\%, leaving
headroom for fragmentation and the real data pipeline). It measures \emph{peak}
memory via \texttt{torch.cuda.max\_memory\_reserved()} (reset before every
probe) rather than current memory, because current memory under-reports the
true high-water mark once autograd has already freed forward-pass activations.
If \texttt{teacher\_fn} is supplied (a distillation run), it is invoked on every
probe batch too, since real training calls the teacher every step and its
memory footprint must be included or the chosen batch size would
out-of-memory the instant real training starts. Returns
\texttt{(micro\_batch, accum\_steps, vram\_fraction)}, where
\texttt{accum\_steps = effective\_batch // micro\_batch} so gradient
accumulation restores the desired effective batch size regardless of what the
hardware can sustain directly. On a CPU-only machine, probing is skipped
entirely (there is no VRAM to probe) and it simply returns
\texttt{min(effective\_batch, max\_probe)} as the micro-batch.

\newpage
\section{pipeline/losses.py}
\role{Defines the two loss functions used across every SegFormer training run: a
combined Dice+BCE loss for direct (non-distilled) training, and a distillation
loss blending ground truth with the teacher's soft probability map.}

\func{class DiceBCELoss(nn.Module)}
\texttt{forward(logits, target)} computes binary cross-entropy directly on raw
logits (numerically stable, versus applying sigmoid then BCE separately), plus
a soft Dice loss (sigmoid probabilities, $2 \cdot |P \cap G| / (|P| + |G| +
\text{smooth})$, with a Laplace-smoothing constant to avoid division by zero on
empty masks), and returns \texttt{bce + (1 - mean\_dice)}. BCE gives stable
gradients everywhere including on empty masks, while Dice directly optimizes the
overlap metric that is ultimately reported --- compensating for BCE's weakness
on a small, class-imbalanced foreground (bruises against a large background).

\func{class DistillSegLoss(nn.Module)}
Implements \texttt{loss = alpha * DiceBCELoss(student, GT) + (1 - alpha) *
BCE(student, teacher\_prob)} --- the same fusion formula used everywhere
distillation appears in this project, so the weighting convention stays
consistent across SegFormer and YOLO distillation rather than being tuned
per-script. \texttt{\_\_init\_\_(alpha=0.75)} stores the mixing weight (ground
truth dominates 3:1 by default) and an internal \texttt{DiceBCELoss} instance
for the supervised term. \texttt{forward(logits, gt, teacher\_prob)} computes the
supervised Dice+BCE term against ground truth, plain BCE-with-logits against the
teacher's (already temperature-scaled, sigmoid-applied) soft probability map,
and returns the alpha-weighted sum.

\newpage
\section{pipeline/mask\_utils.py}
\role{Small mask-related I/O and conversion helpers used by the YOLO training/
evaluation scripts and by the teacher soft-labeling code that feeds YOLO
distillation.}

\func{link\_or\_copy(src: Path, dst: Path) -> None}
Creates \texttt{dst}'s parent directories, then tries to symlink \texttt{dst} to
\texttt{src}; if symlinking fails (e.g. lacking privilege on Windows), falls
back to a real file copy. This lets the pipeline build YOLO-format dataset
directories cheaply via symlinks where supported, without duplicating large
image files on disk, while still working (just using more disk space) where
symlinks are not available. A no-op if \texttt{dst} already exists, so it is
safe to call repeatedly across re-runs.

\func{read\_gt\_mask(path: str) -> np.ndarray}
Reads a grayscale mask via OpenCV, raising \texttt{RuntimeError} naming the path
if unreadable (rather than a confusing downstream \texttt{None} crash), and
binarizes as \texttt{mask > 0} --- robust to masks stored with different
foreground pixel values.

\func{teacher\_prob\_for\_image(model, temperature: float, img\_bgr: np.ndarray, img\_h: int, img\_w: int, device: torch.device) -> np.ndarray}
Generates a teacher soft-probability map for one full-resolution clinical photo:
converts BGR to RGB, \emph{downsamples} to the teacher's trained resolution
(640x640), normalizes with ImageNet mean/std, runs the teacher under
\texttt{no\_grad}, applies temperature-scaled sigmoid, then \emph{upsamples the
resulting probability map} back to the photo's native resolution (clipped to
$[0,1]$ against interpolation overshoot). Downsampling before inference and
upsampling only the result --- rather than running the teacher at native
resolution --- is necessary because SegFormer's self-attention memory scales
with image area, and running at full clinical-photo resolution would
out-of-memory.

\func{save\_semantic\_class\_mask(mask01: np.ndarray, path: Path) -> None}
Writes a binary mask (0 = background, 1 = bruise) as an 8-bit image, matching
Ultralytics' semantic-segmentation label format. Raises \texttt{RuntimeError} on
write failure so a missing/corrupt label file is caught immediately rather than
failing mysteriously during YOLO training.

\func{yolo\_sem\_pred\_mask(result, shape: tuple[int, int]) -> np.ndarray}
Extracts a binary prediction mask from an Ultralytics semantic-segmentation
result object. Ultralytics' semantic task exposes neither \texttt{.masks}
(instance segmentation) nor \texttt{.probs} (classification) --- the dense
per-pixel class map instead lives at \texttt{result.semantic\_mask.data}. This
function pulls that tensor, thresholds it to isolate the bruise class, and
resizes (nearest-neighbor, to avoid blended fractional class values) to the
requested output shape. Returns an all-zero mask if the expected attribute is
missing, so a batch evaluation loop can continue past an unexpected result
rather than crashing.

\newpage
\section{pipeline/metrics.py}
\role{The \textbf{original} pixel-level segmentation metrics module (Dice, IoU,
precision/recall, per-image and aggregate summaries), used by scripts 06--09.
Deliberately left unmodified even after a known correctness issue was found (see
\texttt{metrics\_extended.py}) specifically so that the numbers those scripts
already produced and recorded remain reproducible and comparable to themselves
over time.}

\func{dice\_np(pred: np.ndarray, gt: np.ndarray) -> float}
Computes $2 \cdot |P \cap G| / (|P| + |G|)$, returning $1.0$ when the
denominator is zero (both masks empty). Note: this returns $1.0$ even when
\texttt{pred} is non-empty and \texttt{gt} is empty --- a false positive on a
clean image --- which \texttt{metrics\_extended.py}'s \texttt{dice\_safe}
identifies as a bug and fixes; this file is kept exactly as-is anyway, for the
reproducibility reason stated above.

\func{iou\_np(pred: np.ndarray, gt: np.ndarray) -> float}
Computes $|P \cap G| / |P \cup G|$, returning $1.0$ by convention when the union
is zero (both masks empty, a true negative --- correctly treated as perfect
agreement, unlike the Dice case above).

\func{precision\_recall\_np(pred: np.ndarray, gt: np.ndarray) -> tuple[float, float, int, int, int]}
Computes true/false positives and false negatives via boolean operations.
Precision defaults to $1.0$ if there were no positive predictions at all;
recall defaults to $1.0$ if the ground truth has no positive pixels at all
(both vacuous-truth conventions). Returns all five raw values so callers can
compute dataset-level aggregate precision/recall without recomputing pixel
counts.

\func{compute\_image\_row(pred: np.ndarray, gt: np.ndarray, stem: str) -> dict}
Aggregates Dice, IoU, precision, recall, raw pixel counts, and a
predicted-to-ground-truth positive-pixel ratio (an over/under-segmentation
indicator, set to \texttt{NaN} when ground truth is empty) into one dictionary
row for a single evaluated image, keyed by \texttt{stem} for later joining back
to source images.

\func{summarize(rows: list[dict]) -> dict}
Turns a list of per-image rows into dataset-level aggregates: image count,
mean/median Dice and IoU, counts/rates of zero-Dice images, counts/rates of
``complete miss'' images (model predicted nothing at all where a bruise
existed --- a stricter failure mode than zero Dice alone), mean precision/recall,
and mean predicted-to-ground-truth ratio (with infinite values replaced by
\texttt{NaN} first so they cannot corrupt the mean). Returns an empty dict if
given no rows.

\newpage
\section{pipeline/metrics\_extended.py}
\role{A newer, \textbf{additive} evaluation module used by scripts 09, 10, 11,
and 27, built for a more rigorous benchmark. Its own docstring explains why it
is a separate file rather than a fix to \texttt{metrics.py} in place: scripts
06--09 already produced committed results using the original functions, and
changing them there would silently change already-recorded numbers.}

\func{\_bool\_masks(pred, gt) -> tuple[np.ndarray, np.ndarray]}
Casts both arrays to boolean in one shared place, since every metric function in
this module needs boolean arrays.

\func{\_pixel\_counts(pred\_b, gt\_b) -> tuple[int, int, int, int, int]}
Computes predicted-positive count, ground-truth-positive count, true positives,
false positives, and false negatives in one pass, since Dice/IoU/precision/
recall all need the same five numbers.

\func{dice\_safe(pred, gt) -> float}
The corrected four-case Dice fixing \texttt{metrics.py}'s bug: both masks empty
$\to 1.0$ (a true negative, correctly rewarded); exactly one empty $\to 0.0$
(properly penalizing false positives and complete misses, unlike the original);
both non-empty $\to$ standard Dice.

\func{iou\_safe(pred, gt) -> float}
Same as \texttt{iou\_np} (its empty-mask edge case was not actually broken);
exists mainly for naming symmetry with \texttt{dice\_safe}.

\func{\_boundary(mask: np.ndarray) -> np.ndarray}
Extracts a mask's outer boundary ring by eroding it by one pixel and XOR-ing
against the original --- the standard morphological boundary definition used in
the Surface Dice literature, chosen to match that literature rather than an ad
hoc heuristic.

\func{\_distance\_from\_boundary(mask: np.ndarray) -> np.ndarray}
Computes a per-pixel Euclidean distance transform to the nearest boundary pixel,
in $O(n)$ time for an $n$-pixel image. Returns all-infinity if the mask has no
boundary (i.e. is empty), designed to propagate into downstream \texttt{NaN}/
zero checks rather than raising.

\func{surface\_dice(pred, gt, delta: float = 2.0) -> float}
Implements the boundary-tolerant Surface Dice metric (Nikolov et al. 2018):
what fraction of each predicted/ground-truth boundary lies within \texttt{delta}
pixels of the other boundary. Standard Dice weights every interior pixel
equally, but for bruise segmentation a small boundary misalignment is
clinically acceptable while completely missing the bruise edge is not ---
Surface Dice specifically measures boundary tracking, independent of interior
overlap. Default \texttt{delta=2.0} pixels, per supervisor specification.

\func{asd\_px(pred, gt) -> float}
Average Surface Distance in pixels: the mean boundary-to-boundary distance in
both directions, averaged. Where Surface Dice only answers ``is the boundary
close enough,'' ASD quantifies \emph{how far off} it is on average, a
continuous severity measure Surface Dice alone cannot give.

\func{hd95\_px(pred, gt) -> float}
The 95th-percentile (not the classical 100th-percentile/maximum) Hausdorff
Distance, taken as the max of both directions' 95th percentiles. Using the 95th
percentile rather than the true maximum makes the metric robust to a single
outlier boundary pixel (e.g. one disconnected false-positive fragment) that
would otherwise dominate and make the classical Hausdorff Distance
uninterpretable.

\func{\_safe\_precision(tp, fp) -> float} / \func{\_safe\_recall(tp, fn) -> float}
Same vacuous-truth conventions as \texttt{metrics.py}, explicitly noted to match
scikit-learn's \texttt{zero\_division=1} convention rather than an arbitrary
choice.

\func{compute\_image\_row\_extended(pred, gt, stem, surf\_dice\_delta=2.0) -> dict}
The extended, drop-in replacement for \texttt{compute\_image\_row}: adds the
three boundary metrics above plus explicit \texttt{zero\_dice} and
\texttt{complete\_miss} boolean flags (the latter meaning the model predicted
nothing at all on an image that does contain a bruise).

\func{bootstrap\_ci(values: np.ndarray, n\_bootstrap: int = 2000, ci: float = 0.95, stat: str = "median", seed: int = 0) -> tuple[float, float]}
Computes a bootstrap confidence interval for the median or mean of a metric
column via resampling with replacement and taking empirical percentiles.
Bootstrapping is used instead of a parametric interval because per-image Dice
scores are heavily skewed (many exact zeros from complete misses, many values
near 1.0 from easy cases) rather than normally distributed, which is the
field-standard justification for this method in medical-imaging metrics.
\texttt{NaN} values are dropped before resampling; the seed is fixed for
reproducibility across repeated runs.

\func{wilcoxon\_compare(a: np.ndarray, b: np.ndarray) -> dict}
Runs a paired, two-sided Wilcoxon signed-rank test between two models' per-image
scores on the \emph{same} images, plus a rank-biserial correlation as an effect
size. Wilcoxon is used instead of a t-test because Dice scores are non-normal;
the pairing matters because each row of \texttt{a} and \texttt{b} is the same
test image, making the samples dependent. The effect size is reported alongside
the p-value because with a large test set, p-values alone can flag a trivially
small difference as statistically significant. Returns \texttt{NaN} results
with an explicit warning if fewer than 10 valid (finite, paired) samples remain,
since the test is considered unreliable below that size.

\func{\_aggregate\_column(df: pd.DataFrame, col: str, n\_bootstrap: int) -> dict}
For one metric column: median, IQR, mean, standard deviation, and a 95\%
bootstrap CI. IQR is reported as the statistically correct spread measure for
skewed data; standard deviation is included mainly because readers expect it.

\func{summarize\_extended(rows: list[dict], n\_bootstrap: int = 2000) -> dict}
The extended counterpart to \texttt{metrics.py}'s \texttt{summarize}: aggregates
Dice, IoU, surface Dice, ASD, HD95, precision, and recall via
\texttt{\_aggregate\_column}, plus zero-Dice/complete-miss counts and rates.
Centralizing this in one function guarantees scripts 09, 10, 11, and 27 all
report an identical set of columns, and a new metric added here automatically
propagates to every calling script.

\newpage
\section{pipeline/scheduler.py}
\role{Implements the layer-wise learning-rate and warmup/polynomial-decay
schedule for SegFormer fine-tuning, following the original SegFormer paper's
recipe (Xie et al., 2021). The module's own docstring notes this is a
deliberate departure from a constant-LR convention used elsewhere in the
project for an earlier, unrelated ablation --- this folder tests the paper's
own recipe on its own terms.}

\func{\_is\_no\_decay(name: str, param) -> bool}
Flags parameters that should be excluded from weight decay: names containing
``norm'' or ``bias,'' or any 1-D-or-less tensor (catching bias vectors and norm
scale/shift parameters that might not literally have those substrings in their
name). Applying weight decay to normalization or bias parameters is known to
hurt rather than help training.

\func{build\_param\_groups(model, backbone\_lr: float, head\_lr: float, weight\_decay: float) -> list[dict]}
Splits every trainable parameter into four optimizer groups ---
backbone/decay, backbone/no-decay, head/decay, head/no-decay --- based on
whether a parameter belongs to \texttt{model.backbone} versus
\texttt{model.decode\_head}, and whether \texttt{\_is\_no\_decay} flags it.
Backbone groups get \texttt{backbone\_lr}, head groups get \texttt{head\_lr};
no-decay groups get \texttt{weight\_decay=0.0}. Frozen parameters are skipped
entirely, and empty groups are filtered out (an empty group would error
PyTorch's optimizer). This four-way split is what lets a single
\texttt{AdamW(param\_groups)} call express the paper's layer-wise-LR and
selective-weight-decay recipe.

\func{lr\_multiplier(step: int, total\_steps: int, warmup\_steps: int, power: float = 1.0) -> float}
Returns a pure LR multiplier (not an absolute rate): linear warmup from just
above 0 to 1.0 over the first \texttt{warmup\_steps}, then polynomial decay
(\texttt{power=1.0} is linear, matching this project's use) from 1.0 back to 0
over the remaining steps.

\func{apply\_lr(optimizer, peak\_lrs: list[float], multiplier: float) -> None}
Sets each optimizer parameter group's current learning rate to
\texttt{peak\_lr * multiplier}, iterating the groups in the same order
\texttt{build\_param\_groups} constructed them. Called every training step to
advance the warmup/decay schedule.

\newpage
\section{pipeline/trainer.py}
\role{The shared PyTorch training-and-evaluation engine for all SegFormer runs
--- the B2 teacher, the B0 direct student, and the B0 distilled student. YOLO
training is deliberately handled entirely separately by Ultralytics' own
trainer (scripts 06/08), since YOLO uses its own optimizer, LR schedule, and
data format; this module makes no attempt to unify the two.}

\func{\_load\_temperature(teacher\_dir: Path) -> float}
Reads a calibrated temperature scalar from \texttt{teacher\_dir/temperature.json}
(written by \texttt{02\_calibrate\_teacher.py}). If the file is missing, logs a
warning and returns the safe default $T=1.0$ (no scaling) rather than raising ---
an uncalibrated teacher is still usable, so a missing calibration file is not
treated as fatal.

\func{load\_teacher(teacher\_dir: Path, pretrained: str, device: torch.device, amp: bool = True) -> Callable}
Loads the teacher's best checkpoint (raising \texttt{FileNotFoundError} with an
actionable message if missing) with \texttt{weights\_only=True} (a
security-hardening choice avoiding arbitrary pickle code execution), sets
\texttt{.eval()}, loads its calibrated temperature, and returns a
\texttt{teacher\_fn(x)} closure rather than the raw model --- so the training
loop calls the teacher without knowing anything about its architecture, and a
different teacher implementation could be swapped in later without touching the
training loop. Inside the closure, the forward pass runs under
\texttt{torch.no\_grad()} (the teacher is frozen; the docstring calls this the
single largest memory saving in the distillation setup) and matching AMP
settings, returning temperature-scaled sigmoid probabilities.

\func{evaluate(model: nn.Module, loader: DataLoader, device, threshold: float, amp: bool = True) -> tuple[pd.DataFrame, dict]}
Runs a full no-grad pass over a DataLoader, producing per-image rows (via
\texttt{metrics.compute\_image\_row}) and a dataset summary (via
\texttt{metrics.summarize}). Logits are explicitly cast to fp32 before sigmoid
to avoid overflow near float16's representable range. Training always evaluates
at a fixed 0.50 threshold --- sweeping thresholds every epoch would multiply
evaluation cost roughly 19-fold and add noise to the tracked training curve; the
true best threshold is determined once, after training, in a separate sweep.

\func{\_threshold\_sweep(model, loader, device, thresholds: list[float], amp) -> tuple[pd.DataFrame, float]}
Runs \texttt{evaluate()} once per candidate threshold, returning the full sorted
table and the single best (highest mean Dice) threshold. This sweep must run
against the \textbf{validation} set, never test --- choosing a threshold from
test data would leak test information into model selection and inflate the
reported score.

\func{\_train\_one\_epoch(model, loader, optimizer, sup\_loss, dist\_loss, scaler, accum\_steps, device, clip\_grad, teacher\_fn, total\_steps, warmup\_steps, peak\_lrs, poly\_power, global\_step, desc) -> tuple[float, int]}
Runs one training epoch with gradient accumulation, mixed precision, and the
warmup/decay schedule interleaved. The teacher (if present) is always run
\emph{before} the student under its own no-grad context: this ordering, kept
deliberately consistent with \texttt{find\_optimal\_micro\_batch}'s own probe
ordering, means the teacher's activations are freed before the student's
backward graph is built, roughly halving peak VRAM versus running the student
first. Loss is divided by \texttt{accum\_steps} before backpropagation so
accumulated gradients sum to the same magnitude a full-size batch would
produce; the optimizer step (with gradient clipping) only fires on
accumulation-boundary steps or the epoch's final batch, so a partial trailing
window still gets applied rather than silently dropped.

\func{\_build\_optimizer\_and\_schedule(model, cfg: dict) -> tuple[torch.optim.Optimizer, list[float]]}
Builds the AdamW optimizer from \texttt{scheduler.build\_param\_groups} and
extracts each group's peak learning rate for the warmup/decay schedule.

\func{\_compute\_schedule\_params(n\_batches: int, accum\_steps: int, cfg: dict) -> tuple[int, int, float]}
Computes total optimizer-update steps and warmup steps in terms of actual
gradient-update steps (not raw micro-batches or epochs), since the SegFormer
paper's schedule is inherently defined against optimizer steps and an
epoch-granularity approximation would be imprecise whenever accumulation does
not evenly divide the epoch.

\func{\_save\_run\_config(...) -> None}
Serializes a comprehensive dictionary of hyperparameters and derived quantities
to \texttt{run\_config.json}, so downstream evaluation scripts know exactly how
a given model was trained without re-parsing the original CLI invocation.

\func{\_load\_resume\_state(...) -> tuple} / \func{\_save\_resume\_state(...) -> None}
Save/restore a full training-resume checkpoint (model, optimizer, and AMP
scaler state, plus epoch, best score, patience counter, history, and batch-size
probe results) every epoch, unconditionally --- because GPU sessions can
disconnect mid-training, and writing every epoch bounds lost work to at most
one epoch. This is deliberately decoupled from \texttt{best\_model.pt}: the
resume state tracks ``where training left off,'' while the best-model file
tracks ``the best weights found so far,'' since the most recent epoch may not
be the best one.

\func{train\_pytorch(model, model\_name, run\_dir, train\_df, val\_df, cfg, device, teacher\_fn=None, alpha=0.75) -> dict}
The top-level orchestrator for one full SegFormer run (teacher, direct, or
distilled --- distinguished purely by whether \texttt{teacher\_fn} is
\texttt{None}). Probes or resumes batch sizing, builds DataLoaders and the
optimizer/schedule, writes \texttt{run\_config.json}, then loops epochs calling
\texttt{\_train\_one\_epoch} and \texttt{evaluate}, tracking the best validation
Dice, saving \texttt{best\_model.pt} whenever it improves, writing a resume
checkpoint every epoch regardless, and stopping early once the patience budget
is exhausted. After training, reloads the \emph{best} (not last-epoch) weights,
runs the post-training threshold sweep, and writes \texttt{threshold\_search.csv}
and \texttt{val\_summary.csv}. Returns the full run summary dictionary.

\newpage
\section{pipeline/yolo\_threshold\_temp.py}
\role{Temperature scaling and threshold sweeping for YOLO26n-sem. Ultralytics'
own \texttt{.predict()} converts raw logits straight to a hard class map via
argmax internally, exposing no probability map or threshold control --- this
module bypasses that entirely by calling the underlying \texttt{nn.Module}
directly. For a 2-class semantic head, BCE training pushes logits toward
$\pm\infty$, producing a near-binary probability histogram with almost nothing
near 0.5; temperature scaling (dividing logits by $T>1$ before sigmoid) spreads
that histogram back out so a threshold sweep has something meaningful to
differentiate.}

\func{yolo\_raw\_class\_logits(yolo\_model, x: torch.Tensor, out\_hw: tuple[int, int]) -> torch.Tensor}
Runs YOLO's raw \texttt{nn.Module} directly (bypassing \texttt{.predict()} and
its internal argmax), and resizes the raw \texttt{[B, nc, h, w]} output to
\texttt{out\_hw} if it does not already match --- this is the exact call site
where the 640x640 resize happens inside \texttt{benchmark\_640.py}'s timed
section.

\func{bruise\_prob\_from\_logits(class\_logits: torch.Tensor, temperature: float = 1.0) -> torch.Tensor}
Converts \texttt{[B, nc, H, W]} class logits into a single \texttt{[B, H, W]}
bruise probability map: for a 2-class head, takes the class-logit difference
(bruise minus background) and applies sigmoid --- mathematically identical to
2-class softmax, and directly analogous to SegFormer's single-channel sigmoid.
Dividing by \texttt{temperature} before sigmoid is the de-saturation step;
$T=1.0$ is the model's native (near-binary) output, $T>1.0$ spreads the
histogram back toward 0.5.

\func{evaluate\_yolo\_raw(yolo\_model, loader, device, threshold: float, temperature: float) -> tuple[pd.DataFrame, dict]}
Same shape/contract as \texttt{trainer.evaluate()}, but for YOLO's raw
\texttt{nn.Module} plus this module's own sigmoid/threshold postprocessing
instead of Ultralytics' \texttt{.predict()} plus argmax.

\func{threshold\_temperature\_sweep(yolo\_model, loader, device, thresholds: list[float], temperatures: list[float]) -> pd.DataFrame}
A joint grid sweep over threshold x temperature on validation, scored by mean
Dice --- the same simple grid-sweep method used for SegFormer, just crossed
with an extra temperature axis. Returns the full grid (best-first) so the
empty-gap/de-saturation pattern is visible in the saved CSV, not only the
single winning row.

\func{run\_threshold\_search(weights\_path: str, val\_df, cfg: dict, device, run\_dir: Path) -> tuple[pd.DataFrame, float, float]}
The full validation-side pipeline for one YOLO run: loads \texttt{best.pt},
extracts the raw \texttt{nn.Module}, sweeps temperature x threshold on
validation, saves \texttt{threshold\_search.csv} (same filename convention as
SegFormer, so downstream scripts read both uniformly), and returns the full
grid plus the best threshold and temperature.

\func{load\_yolo\_threshold\_temp(run\_dir: Path) -> tuple[float, float]}
Reads the best validation \texttt{(threshold, temperature)} pair from
\texttt{threshold\_search.csv}, defaulting temperature to 1.0 for any older CSV
written before temperature scaling existed.

\func{load\_yolo\_model(model\_name: str, paths: dict, device: torch.device) -> tuple[nn.Module, float, float, Path]}
Loads a trained YOLO run's raw \texttt{nn.Module} (deep-copied out of the
Ultralytics wrapper) plus its calibrated \texttt{(threshold, temperature)} pair.
Returning the raw module (not the wrapper) means a caller can never accidentally
fall back to Ultralytics' own argmax postprocessing.

\newpage
\section{pipeline/benchmark\_640.py}
\role{A device-agnostic (CPU or CUDA) controlled 640x640 deployment-latency
benchmark. This exists because \texttt{scripts/29\_benchmark\_inference\_all\_models.py}
raises immediately without CUDA --- reasonable for final GPU numbers, but it
means the benchmark cannot be exercised at all while the cluster is down. Every
timing helper here is device-agnostic: synchronization is a no-op off CUDA,
peak memory is reported as \texttt{NaN} off CUDA, and fp16 is skipped (not
attempted) off CUDA. It also fixes a real inconsistency in script 29: YOLO is
preprocessed with the exact same ImageNet-normalized transform
(\texttt{pipeline.data.get\_augmentation}) as SegFormer, matching how YOLO's
threshold and temperature were actually calibrated
(\texttt{yolo\_threshold\_temp.run\_threshold\_search} and the script that
produced the project's reported test-set Dice both go through
\texttt{BruiseDataset}, which normalizes with ImageNet mean/std) --- script 29's
own YOLO preprocessing instead used a plain 0--1 pixel scale, a different input
distribution than the one its own loaded threshold was calibrated against.}

\func{class SourceImage (frozen dataclass): stem: str, path: str}
A benchmark image \emph{reference} only --- no decoded pixels. Deliberately not
preloaded: each fixed-test image is roughly 4022x6024x3 pixels ($\approx$72MB
decoded), and holding all 185 of them in memory at once needs about 13GB, which
actually crashed with an out-of-memory error on this laptop mid-preload.

\func{read\_source\_image(source: SourceImage) -> np.ndarray}
Decodes one image (OpenCV read, BGR to RGB) on demand, called fresh immediately
before each timed inference and never cached --- keeping peak memory to roughly
one image at a time while keeping disk decode excluded from the timed
measurement, exactly as before the memory fix.

\func{class BenchmarkSummary (dataclass)}
The schema of one row of the output summary CSV: model identity, checkpoint,
device, precision, threshold, temperature, parameter counts, image/repeat
counts, mean/median/p90/p95/std latency, FPS from mean and from median, peak GPU
memory, and a fixed description of exactly what was timed.

\func{class Fixed640Adapter}
The common base class for the controlled benchmark.
\texttt{\_\_init\_\_(...)} stores the model, threshold, temperature, device, and
precision, and builds \texttt{self.transform} via
\texttt{get\_augmentation(training=False, ...)} --- the identical transform
\texttt{BruiseDataset} uses at evaluation time. The \texttt{autocast\_enabled}
property is \texttt{True} only for fp16 on a CUDA device (there is no fp16
speed benefit to measure on CPU in this codebase).
\texttt{parameter\_count} returns \texttt{count\_params(self.model)[0]}.
\texttt{\_preprocess(image\_rgb)} applies the transform, adds a batch dimension,
and moves the tensor to the device. \texttt{infer\_full(image\_rgb)} is abstract:
it must take an RGB image already in RAM and return a CPU \texttt{uint8} binary
mask of shape $640 \times 640$; subclasses implement only the one line that
differs between architectures.

\func{class SegFormer640Adapter(Fixed640Adapter)}
\texttt{infer\_full} preprocesses, runs the model under autocast (no extra
resize needed --- \texttt{SegformerWrapper} already upsamples logits to input
resolution internally), divides by temperature (always 1.0 for SegFormer in
this project) before sigmoid, thresholds, and returns the CPU \texttt{uint8}
mask.

\func{class YOLOSemantic640Adapter(Fixed640Adapter)}
\texttt{infer\_full} preprocesses through the \emph{same} ImageNet-normalized
transform as SegFormer, calls \texttt{yolo\_raw\_class\_logits} (never
Ultralytics' \texttt{.predict()}), applies \texttt{bruise\_prob\_from\_logits}
with the calibrated temperature, thresholds, and returns the mask.

\func{build\_adapter(spec: dict, paths: dict, device, precision: str) -> Fixed640Adapter}
Builds one adapter from a \texttt{configs/benchmark\_640\_models.yaml} entry,
dispatching on family and calling \texttt{load\_segformer\_model} or
\texttt{load\_yolo\_model} --- deliberately reusing the exact same
checkpoint/threshold resolution logic script 29 uses, so ``which checkpoint''
and ``which threshold'' can never drift between the two benchmark scripts.
Raises \texttt{ValueError} for an unrecognized family.

\func{load\_fixed\_test\_images(paths: dict, max\_images: int | None) -> list[SourceImage]}
Lists the fixed test-set images to benchmark (stems and paths only, per
\texttt{SourceImage}'s memory rationale above). Raises \texttt{RuntimeError} if
the manifest yields zero images.

\func{\_synchronize(device: torch.device) -> None}
A no-op unless the device is CUDA, in which case it calls
\texttt{torch.cuda.synchronize}. GPU work is dispatched asynchronously, so
without this, a timer would measure kernel-launch latency rather than actual
compute time; CPU execution is already synchronous, so nothing is needed there.

\func{\_summarize(values\_ms) -> dict[str, float]}
Computes mean, median, p90, p95, standard deviation, and FPS-from-mean/
FPS-from-median from a list of per-call latencies. Raises \texttt{ValueError} on
an empty input.

\func{benchmark\_adapter(adapter, images, *, warmup: int, repeats: int) -> tuple[BenchmarkSummary, pd.DataFrame]}
The core timing loop. Refuses outright (raises \texttt{ValueError}) if fp16 is
requested off a CUDA device, rather than silently timing fp32 under an
``fp16'' label. Runs \texttt{warmup} untimed iterations first (absorbing
first-call costs like lazy kernel compilation), asserting every returned mask
is exactly $640{\times}640$ \texttt{uint8} or raising immediately. Then, for
\texttt{repeats} passes over every image: decodes the image fresh (before the
timer starts), synchronizes, starts the timer, calls
\texttt{adapter.infer\_full}, synchronizes again, and records the elapsed
milliseconds, re-checking shape/dtype every single iteration. Finally computes
\texttt{\_summarize} and peak GPU memory (\texttt{NaN} off CUDA), and assembles
the \texttt{BenchmarkSummary}.

\func{run\_benchmark\_640(*, paths\_config, common\_config, models\_config, output\_dir, device\_name, precisions, max\_images, warmup, repeats) -> pd.DataFrame}
The top-level orchestrator called by \texttt{scripts/31\_benchmark\_640\_local.py}.
Loads and validates configs; refuses to run if \texttt{common\_train.yaml}'s
image size is not $640{\times}640$ (this module is hardcoded to the 640
protocol, and would rather fail loudly than silently benchmark at the wrong
resolution); resolves the device (raising if CUDA was requested but
unavailable); loads the benchmark images once. For every requested precision
and every enabled model in the registry, it builds an adapter and benchmarks
it, catching any exception broadly (not just a missing file) so one model
failing to load --- a missing run folder, or an incompatible checkpoint, per
\texttt{\_pick\_matching\_checkpoint} above --- prints a \texttt{SKIP} message
and moves on instead of aborting the entire run, matching script 29's own
SKIP behavior. After every model, memory is explicitly freed before the next
one loads. Finally, writes the summary CSV/JSON, per-image latency CSV, and a
metadata JSON recording every config path and setting used for that run.

Module constant: \texttt{IMAGE\_SIZE = 640}.

\newpage
\part{Pipeline Scripts (\texttt{scripts/})}

\section{scripts/00\_build\_split.py}
\role{The pipeline's foundational stage: builds a fresh, leakage-checked
subject-level train/val split, excluding any subject that appears in the fixed
test manifest. Splitting at the \emph{subject} level (not image level) matters
because one forensic subject can contribute multiple images (different
lighting/angle) --- an image-level split could put the same person's images in
both train and val, leaking skin-tone and bruise-appearance information and
inflating reported Dice. Building the split fresh here (rather than reusing one
from another experiment) keeps every pipeline run self-contained and avoids
silently importing another experiment's seed or contamination.}

\func{main()}
Reads and normalizes both the train and fixed-test manifests, computes the
candidate subject pool (train subjects minus test subjects), shuffles with a
seeded RNG, splits off \texttt{val-fraction} of subjects into validation
(guaranteeing at least one), and performs an explicit three-way leakage check
(train/val, train/test, val/test), raising \texttt{RuntimeError} immediately if
any overlap is found --- treated as fatal since silent leakage would invalidate
every downstream Dice number. Writes \texttt{train\_val\_split.csv} and
\texttt{split\_summary.csv}.

\begin{itemize}[nosep]
\cliarg{--paths}{default \texttt{configs/paths.yaml}}
\cliarg{--val-fraction}{float, default 0.18}
\cliarg{--seed}{int, default 42}
\end{itemize}

\section{scripts/01\_train\_segformer\_b2\_teacher.py}
\role{Trains the SegFormer-B2 teacher (27.3M parameters) from an
ImageNet-pretrained backbone with a randomly initialized binary decode head,
using the paper's layer-wise LR recipe. Must run before any distillation step,
since those steps consume this teacher's soft labels --- a weaker teacher caps
the ceiling of every student.}

\func{main()}
Skips training if \texttt{best\_model.pt} already exists (unless
\texttt{--force-retrain}). Loads the train/val split, builds the model, and
delegates the entire training loop to \texttt{pipeline.trainer.train\_pytorch}
with \texttt{teacher\_fn=None} (this model has no teacher of its own).

\begin{itemize}[nosep]
\cliarg{--paths, --common}{config paths}
\cliarg{--force-retrain}{flag; re-train even if a checkpoint exists}
\end{itemize}

\section{scripts/02\_calibrate\_teacher.py}
\role{Applies temperature scaling (Guo et al. 2017) to the trained teacher's
logits, fitting a single scalar temperature via L-BFGS on validation-set
logits. Necessary because BCE-trained logits saturate toward $\pm\infty$, making
soft labels nearly indistinguishable from hard ground truth and defeating the
purpose of distillation; dividing by $T>1$ spreads the distribution so students
can see the teacher's real uncertainty near decision boundaries.}

\func{\_collect\_val\_logits(model, loader, device) -> tuple[torch.Tensor, torch.Tensor]}
Runs one no-grad pass over the full validation loader, concatenating all logits
and targets. Collected all at once (not streamed per L-BFGS iteration) because
L-BFGS needs a stable full-dataset loss/gradient at every step, and the
validation set is small enough that this is cheap and done only once.

\func{\_calibrate(logits, targets) -> float}
Optimizes a scalar \texttt{log\_t} via \texttt{torch.optim.LBFGS}, minimizing
BCE-NLL of \texttt{logits/T} against targets; optimizing in log-space
constrains $T>0$ automatically. Returns $T = \exp(\text{log\_t})$.

\func{main()}
Skips if \texttt{temperature.json} already exists; raises if the teacher
checkpoint is missing. Collects validation logits, computes NLL before and
after calibration, writes both to \texttt{temperature.json}, and logs a warning
if the fitted $T < 1.0$ (unusual --- would imply an under-confident model,
atypical for BCE training).

\begin{itemize}[nosep]
\cliarg{--paths, --common}{config paths}
\end{itemize}

\section{scripts/03\_train\_segformer\_b0\_direct.py}
\role{Trains the small SegFormer-B0 student (3.7M parameters) directly on
ground truth only --- no teacher, no distillation --- using identical
hyperparameters to the teacher. This is the baseline Track A's apples-to-apples
comparison measures distillation (script 05) against: since both use identical
recipes, any Dice improvement can be attributed specifically to the teacher
signal, not a confounded hyperparameter difference.}

\func{main()}
Structurally identical to script 01's \texttt{main()}: builds the B0 model and
calls \texttt{train\_pytorch(..., teacher\_fn=None)}.

\begin{itemize}[nosep]
\cliarg{--paths, --common}{config paths}
\cliarg{--force-retrain}{flag}
\end{itemize}

\section{scripts/04\_optuna\_alpha\_segformer\_b0.py}
\role{Runs a TPE (Tree-structured Parzen Estimator) Bayesian search via Optuna
to find the optimal alpha (the ground-truth-versus-teacher weighting) for
SegFormer-B0 distillation. TPE is used instead of a grid search because it
concentrates sampling in the high-Dice region and typically converges faster.
Each trial trains only a short number of epochs, since the goal is
\emph{relative ranking} of alpha values, and rankings stabilize well before full
convergence.}

\func{\_run\_trial(trial, train\_df, val\_df, search\_cfg, paths, device, teacher\_fn, out\_dir) -> float}
Suggests an alpha from the search space, trains a \textbf{fresh} model for every
trial (avoiding contamination from a previous alpha's partially-adapted
weights), and returns the resulting mean Dice. Returns $0.0$ on any exception
(logged) so one failed trial does not abort the whole study; frees GPU memory
after every trial.

\func{main()}
Overrides epochs/patience to a short search-specific value (with patience equal
to epochs, so early stopping cannot bias short trials), creates an Optuna study
with a seeded TPE sampler and SQLite storage (resumable if interrupted), runs
the configured number of trials, and writes the best alpha (falling back to a
configured default only if Optuna's result is somehow missing the key) plus the
full trial history.

\begin{itemize}[nosep]
\cliarg{--paths, --common}{config paths (search bounds/trial count come from \texttt{common\_train.yaml})}
\end{itemize}

\section{scripts/05\_train\_segformer\_b0\_distilled.py}
\role{Trains the final, full-length distilled SegFormer-B0 checkpoint using
knowledge distillation from the calibrated teacher at the Optuna-found alpha.
Feeds directly into Track A/B evaluation.}

\func{\_load\_optuna\_alpha(project\_root, cfg) -> float}
Reads the Optuna best-alpha CSV, raising \texttt{RuntimeError} (never silently
defaulting) if it is missing --- a silent default would misrepresent results as
using an Optuna-optimized alpha when they did not.

\func{main()}
Resolves alpha from a CLI override or the Optuna result, loads the calibrated
teacher, builds a fresh B0 model, and calls
\texttt{train\_pytorch(..., teacher\_fn=teacher\_fn, alpha=alpha)}.

\begin{itemize}[nosep]
\cliarg{--paths, --common}{config paths}
\cliarg{--alpha}{float, optional override of the Optuna-found value}
\cliarg{--force-retrain}{flag}
\end{itemize}

\section{scripts/06\_train\_yolo\_sem\_direct.py}
\role{Trains YOLO26n-sem directly (no distillation) using Ultralytics' own
training pipeline rather than the custom SegFormer trainer --- a deliberate
choice since Ultralytics tightly couples its backbone/neck/head training with
its own AMP, EMA, and augmentation machinery, and extracting the raw
\texttt{nn.Module} for a custom loop would risk losing that stability.}

\func{build\_split(df, out\_dir, split) -> None}
Builds the Ultralytics-format dataset for one split: symlinks each source image
(avoiding full copies of large forensic images) and converts each binary
ground-truth mask into a class-index PNG, since Ultralytics' semantic task
expects class-index masks, not confidence floats.

\func{\_write\_data\_yaml(run\_dir, data\_dir, train\_n, val\_n) -> Path}
Writes the Ultralytics \texttt{data.yaml} descriptor Ultralytics' training API
requires, which also serves as a permanent record of exactly what data trained
the model.

\func{main()}
Builds the dataset (skipping rebuild if already complete), then trains (or
resumes, or skips if already trained) via Ultralytics' own \texttt{.train()}
call with \texttt{task="semantic"}, 640 image size, cosine LR, and mosaic
augmentation disabled for the final epochs (so the model fine-tunes on
single-image crops matching test-time input). Afterward evaluates on
validation using Ultralytics' native \texttt{.predict()} (argmax, no
threshold/temperature tuning) and writes the summary CSVs.

\begin{itemize}[nosep]
\cliarg{--paths, --common}{config paths}
\cliarg{--force-rebuild-dataset, --force-retrain}{flags}
\end{itemize}

\section{scripts/07\_optuna\_alpha\_yolo\_sem.py}
\role{The analogous Optuna alpha search for YOLO, but using \emph{offline}
knowledge distillation: Ultralytics' training loop has no hook for injecting
custom per-batch soft targets, so fused pseudo-masks (ground truth blended with
the temperature-scaled teacher probability, weighted by alpha and thresholded)
must be precomputed to disk before each trial trains.}

\func{build\_split(df, out\_dir, split, alpha, teacher, temperature, cfg, device) -> None}
For the training split with a teacher present: computes the teacher's
probability map per image, fuses it with ground truth by alpha, and thresholds
the fused mask. For the validation split, always uses pure (unfused) ground
truth, since validation masks are the scoring reference and must not be
softened.

\func{run\_trial(alpha, trial\_number, weights\_path, out\_dir, train\_df, val\_df, cfg, teacher, temperature, device) -> float}
Builds a fresh pseudo-mask dataset per trial, trains a short YOLO run (early
stopping disabled, warmup capped, logging suppressed), evaluates on validation
via Ultralytics' native prediction, and returns mean Dice. Deletes the trial's
dataset afterward (retaining only the Ultralytics run logs) to avoid exhausting
disk space across many trials.

\func{main()}
Loads the raw teacher \texttt{nn.Module} directly (not via the
\texttt{teacher\_fn} wrapper, since the pseudo-mask builder needs to run its own
preprocessing on raw images), creates a seeded TPE Optuna study, runs the
configured trial count, and writes the best-alpha CSV and trial history.

\begin{itemize}[nosep]
\cliarg{--paths, --common}{config paths}
\end{itemize}

\section{scripts/07b\_threshold\_yolo\_direct.py}
\role{Performs the joint (temperature, threshold) grid sweep on validation for
the \emph{direct} YOLO model, bypassing Ultralytics' native prediction path
entirely. All sweep logic lives in the shared
\texttt{pipeline/yolo\_threshold\_temp.py} helper, not duplicated here.}

\func{main()}
Raises if the direct YOLO weights are missing; skips if
\texttt{threshold\_search.csv} already exists (unless \texttt{--force-rerun});
calls \texttt{run\_threshold\_search} and logs the best pair plus the top rows
of the grid.

\begin{itemize}[nosep]
\cliarg{--paths, --common}{config paths}
\cliarg{--force-rerun}{flag}
\end{itemize}

\section{scripts/07c\_threshold\_yolo\_distilled.py}
\role{Identical logic to 07b, applied to the \emph{distilled} YOLO checkpoint.
Kept as a separate script (rather than reusing 07b's result) because
distillation changes the model's learned decision boundary and logit scale, so
the jointly optimal (temperature, threshold) pair typically differs from the
direct model's even with an identical architecture.}

\func{main()}
Structurally identical to 07b's, pointed at the distilled run directory.

\begin{itemize}[nosep]
\cliarg{--paths, --common}{config paths}
\cliarg{--force-rerun}{flag}
\end{itemize}

\section{scripts/08\_train\_yolo\_sem\_distilled.py}
\role{Trains the final, full-length distilled YOLO checkpoint using the
Optuna-found alpha from script 07, following the same offline pseudo-mask
distillation approach.}

\func{\_load\_optuna\_alpha(project\_root) -> float}
Reads the YOLO Optuna best-alpha CSV, raising if missing (mirroring script 05's
rationale --- no silent default).

\func{build\_split(df, out\_dir, split, alpha, teacher, temperature, cfg, device) -> None}
Same fusion logic as script 07's version, wrapped in a progress bar since
building pseudo-masks requires one GPU forward pass per training image.

\func{main()}
Resolves alpha, builds/rebuilds the pseudo-mask dataset (freeing the teacher
from GPU memory before YOLO training starts), trains via Ultralytics with the
same recipe as script 06, evaluates on validation with native prediction, and
records alpha in \texttt{run\_config.json}.

\begin{itemize}[nosep]
\cliarg{--paths, --common}{config paths}
\cliarg{--alpha}{float, optional override}
\cliarg{--force-rebuild-dataset, --force-retrain}{flags}
\end{itemize}

\section{scripts/09\_evaluate\_test.py}
\role{A fast ``sanity check'' evaluation of all 5 models on the fixed held-out
test set --- basic Dice/IoU/Precision/Recall and simple FPS timing, no
bootstrapping or statistical tests, distinct from the rigorous Track A/B
scripts. YOLO here uses Ultralytics' native argmax prediction (no
threshold/temperature), contrasted against script 09b's temperature-scaled
version.}

\func{\_load\_best\_threshold(run\_dir) -> float}
Reads a SegFormer run's validation-selected threshold; raises rather than
defaulting to 0.5, so results stay comparable to the Track A/B scripts, which
always use the validation-selected threshold.

\func{eval\_pytorch(run\_name, pretrained, run\_dir, test\_df, cfg, device, out\_dir) -> dict}
Loads the model and threshold, evaluates the fixed test set at batch size 1
(chosen so per-image timing is accurate rather than amortized across a batch),
timing each forward pass with GPU synchronization, and aggregates per-image
metrics.

\func{eval\_yolo(run\_name, run\_dir, test\_df, cfg, out\_dir) -> dict}
Evaluates YOLO via Ultralytics' native \texttt{.predict()} (argmax, no
threshold), timing with wall-clock \texttt{time.perf\_counter()} since
\texttt{.predict()} is already a blocking call that includes any GPU wait.

\func{main()}
Iterates the 3 SegFormer runs (skipping any without a checkpoint) calling
\texttt{eval\_pytorch}, then the 2 YOLO runs calling \texttt{eval\_yolo}, and
writes a combined, Dice-sorted comparison CSV.

\begin{itemize}[nosep]
\cliarg{--paths, --common}{config paths}
\cliarg{--force}{flag; re-evaluate even if outputs exist}
\end{itemize}

\section{scripts/09b\_evaluate\_yolo\_temp\_scaled.py}
\role{Re-evaluates the two YOLO models on the fixed test set using the
validation-selected (temperature, threshold) pairs, since Ultralytics'
\texttt{.predict()} pipeline hard-codes an internal argmax step that cannot be
intercepted to inject temperature scaling. Outputs are written to a parallel
directory so both this script's and script 09's results are preserved for
comparison.}

\func{\_load\_temp\_threshold(run\_dir) -> tuple[float, float]}
Reads the validation-selected \texttt{(threshold, temperature)} pair, raising if
the sweep has not been run --- falling back to $(T{=}1, \text{threshold}{=}0.5)$
would be misleading, since YOLO's near-binary logits specifically need
$T \gg 1$ for a threshold to be meaningful.

\func{eval\_yolo\_with\_temp\_threshold(run\_name, run\_dir, test\_df, cfg, device, out\_dir) -> dict}
Deep-copies the raw \texttt{nn.Module} out of the Ultralytics wrapper (to avoid
mutating the wrapper's internal state), evaluates using
\texttt{yolo\_raw\_class\_logits} and temperature-scaled \texttt{bruise\_prob\_from\_logits}
directly, bypassing Ultralytics' postprocessing.

\func{main()}
For each YOLO run, skips (with a warning, since this script is supplementary)
if the threshold sweep has not been run; otherwise evaluates and writes a
combined comparison CSV.

\begin{itemize}[nosep]
\cliarg{--paths, --common}{config paths}
\cliarg{--force}{flag}
\end{itemize}

\section{scripts/10\_consolidate\_results.py}
\role{A lightweight, no-GPU-work aggregation script merging each model's
validation and (if available) test summaries into one wide CSV alongside
training-recipe metadata --- purely for reporting convenience, distinct from
the deeper statistical Track A/B outputs.}

\func{main()}
For each of the 5 fixed run names, reads \texttt{val\_summary.csv} (skipping
with a message if absent), splits columns into identity/recipe fields versus
metric fields, optionally merges in \texttt{test\_summary.csv}, and writes the
combined table.

\begin{itemize}[nosep]
\cliarg{--paths}{config path}
\end{itemize}

\section{scripts/10\_track\_a\_evaluate.py}
\role{``Track A'': a strict, apples-to-apples benchmark of only the three
SegFormer variants trained under identical conditions, styled after Gut et
al. 2022. YOLO is deliberately excluded here because it uses a fundamentally
different optimizer, schedule, loss, and data format --- including it would
violate the ``identical conditions'' premise (that comparison is Track B's
job). Reports Dice/IoU/Precision/Recall, boundary metrics, fairness breakdowns,
three categories of inference speed, and statistical rigor via bootstrap
confidence intervals and Wilcoxon tests against the B0-direct baseline.}

\func{load\_best\_threshold(run\_dir) -> float}
Same pattern as other scripts --- reads the validation-selected threshold,
raising rather than silently defaulting.

\func{\_gpu\_timed(fn, n\_warmup, n\_iters) -> tuple[float, float, float]}
A generic GPU timing helper: warms up untimed, then times \texttt{n\_iters}
calls with \texttt{cuda.synchronize()} bracketing each, since GPU work is
asynchronous and an unsynchronized timer would measure kernel-launch latency
rather than actual compute time.

\func{\_build\_single\_image\_tensor(image\_path, cfg, device) -> torch.Tensor}
Loads one image through \texttt{BruiseDataset} and returns a GPU-ready tensor,
built outside the timed function since the controlled benchmark measures only
the forward pass and postprocessing, deliberately excluding disk I/O (which is
not reproducible due to OS disk-cache effects).

\func{benchmark\_speed(model, best\_thr, benchmark\_img, test\_paths, cfg, device, n\_warmup, n\_iters) -> dict}
Measures three timing categories: \texttt{raw\_forward} (bare model call,
isolating architecture speed), \texttt{mask\_output} (forward plus sigmoid plus
threshold), and \texttt{end\_to\_end} (full realistic pipeline including disk
read, preprocessing, GPU transfer, and threshold), the last measured once per
real test image to average out disk I/O variance rather than repeating a
single cached image.

\func{\_run\_inference\_loop(model, loader, best\_thr, device, amp, surf\_dice\_delta) -> list[dict]}
One pass over the test loader computing extended per-image metrics (including
surface Dice) for every image, keeping per-image rows (not just an aggregate)
so later analyses --- per-skin-tone breakdowns, Wilcoxon tests, outlier
inspection --- do not require re-running inference.

\func{eval\_track\_a(run\_name, pretrained, run\_dir, test\_df, cfg, device, out\_dir, benchmark\_img, test\_paths, surf\_dice\_delta=2.0, n\_bootstrap=2000, n\_warmup=15, n\_iters=100) -> dict}
The full per-model pipeline: loads the validation threshold, loads
\texttt{best\_model.pt} (never last-epoch weights, avoiding an overfit
checkpoint), runs inference, saves per-image results \emph{before} aggregation
(so a crash during aggregation cannot lose raw predictions), computes the
extended summary, runs the speed benchmark separately, and returns everything
plus the per-image Dice array (used later for Wilcoxon tests, not persisted to
CSV).

\func{\_run\_wilcoxon\_tests(all\_results, baseline\_name, out\_dir) -> list[dict]}
Runs a Wilcoxon test between each non-baseline model's per-image Dice and the
baseline's (\texttt{segformer\_b0\_direct}, chosen as the simplest student with
no teacher --- directly answering ``does distillation help?'').

\func{main()}
Iterates the 3 SegFormer runs (reusing cached results if already evaluated
unless \texttt{--force}), runs the Wilcoxon tests against the baseline, and
writes a Dice-sorted comparison CSV plus the Wilcoxon table.

\begin{itemize}[nosep]
\cliarg{--paths, --common}{config paths}
\cliarg{--n-warmup, --n-iters}{GPU timing warmup/timed iteration counts (default 15/100)}
\cliarg{--n-bootstrap}{bootstrap resamples for CI (default 2000)}
\cliarg{--surf-delta}{Surface Dice tolerance in pixels (default 2.0)}
\cliarg{--force}{flag}
\end{itemize}

\section{scripts/11\_track\_b\_evaluate.py}
\role{``Track B'': a deployment-realistic comparison of all 5 models, each
using its own best/realistic recipe rather than a forced-identical protocol.
The docstring is explicit that Track B can only support ``better in our
realistic setup,'' not ``architecturally superior'' --- that stronger claim
requires Track A's controlled comparison, since SegFormer and YOLO use
fundamentally different optimizers/losses/preprocessing here.}

\func{load\_segformer\_threshold(run\_dir) -> float} / \func{load\_yolo\_threshold\_temp(run\_dir) -> tuple[float, float]}
Same validation-selected threshold (and, for YOLO, temperature) loading as
Track A, extended to cover both model families.

\func{\_gpu\_timed(fn, n\_warmup, n\_iters) -> tuple[float, float, float]}
Identical helper to Track A's, for the identical reason (avoiding measurement
of async kernel-launch latency instead of real compute time).

\func{\_e2e\_times\_segformer(model, best\_thr, test\_paths, cfg, device) -> np.ndarray}
Measures true end-to-end wall-clock latency for SegFormer across every test
image: disk read, color conversion, preprocessing, GPU transfer, forward,
sigmoid, threshold, and device-to-host transfer, each step timed together.

\func{\_e2e\_times\_yolo(yolo\_nn, best\_thr, best\_temp, test\_paths, cfg, device) -> np.ndarray}
The YOLO analogue, using \texttt{yolo\_raw\_class\_logits} and
temperature-scaled \texttt{bruise\_prob\_from\_logits} in place of a plain
sigmoid.

\func{eval\_segformer\_track\_b(...) -> dict}
Nearly identical to Track A's per-model evaluation, but written under
\texttt{track\_b\_evaluation/} and additionally recording parameter count and
\texttt{temperature\_used=1.0}, so the final comparison table can show
parameter counts side by side across SegFormer and YOLO.

\func{eval\_yolo\_track\_b(...) -> dict}
The YOLO analogue: deep-copies the raw \texttt{nn.Module}, evaluates with
temperature-scaled logits, and runs the same three-category speed benchmark
using YOLO-specific logit extraction in place of a plain forward call.

\func{\_print\_result(run\_name, res) -> None}
A logging helper printing one formatted summary line per model: median Dice
with its 95\% CI, HD95, complete-miss rate, and both raw-forward and
end-to-end FPS.

\func{main()}
Iterates the 3 SegFormer runs then the 2 YOLO runs (skipping any missing
weights or threshold sweep), runs Wilcoxon tests against the B0-direct
baseline for every other model, and writes the combined comparison and
Wilcoxon tables.

\begin{itemize}[nosep]
\cliarg{--paths, --common}{config paths}
\cliarg{--n-warmup, --n-iters, --n-bootstrap, --surf-delta, --force}{same meaning as Track A}
\end{itemize}

\section{scripts/24\_consolidate\_all\_results.py}
\role{A final ``gather everything'' step copying every human-readable results
artifact from scripts 00--23 (plus ad-hoc benchmark/audit scripts) into one
self-contained mirrored folder, so results can be reviewed or zipped without
hunting across many scattered directories. Explicitly distinguished from
\texttt{10\_consolidate\_results.py}, which merges/transforms numbers into one
CSV; this script only copies files, never transforms anything. Model
checkpoint weights are never copied (they would make the folder multi-GB,
defeating the point), and a missing source directory is logged as a warning in
the output manifest rather than raised as an error, since the script must be
runnable at any point in the pipeline's lifetime.}

\func{\_copy\_filtered\_tree(src\_dir, dst\_dir, manifest: list[str]) -> None}
Recursively copies only files with an allowed extension (CSV, JSON, TXT, XLSX,
LOG) from \texttt{src\_dir} into \texttt{dst\_dir}, preserving relative
structure; anything else (weights, images, caches) is simply skipped by the
extension filter with no special-casing needed.

\func{\_copy\_file(src, dst, manifest) -> None}
Copies a single named file, recording a MISSING line in the manifest instead of
raising if the source does not exist.

\func{main()}
Copies the split directory, every per-model run directory's result files, the
Optuna search results, every evaluation folder, and the ITA labels directory
(if configured) into the output tree, then writes the accumulated
\texttt{MANIFEST.txt}.

\begin{itemize}[nosep]
\cliarg{--paths}{config path}
\cliarg{--out-name}{default \texttt{CONSOLIDATED\_RESULTS}}
\end{itemize}

\section{scripts/26\_migrate\_legacy\_segformer\_checkpoints.py}
\role{A one-time migration utility rewriting SegFormer checkpoint parameter
names between two different internal layouts HuggingFace \texttt{transformers}
has used across versions --- discovered 2026-07-07 while running scripts 25 and
12, and confirmed identical across two independently trained checkpoints.
Rewriting the checkpoint file's own keys (rather than patching every loading
call site in \texttt{pipeline/trainer.py} and every evaluation script) fixes
every caller at once without touching a single \texttt{.py} file --- and this
project's own ground rule forbids modifying \texttt{pipeline/} or scripts
00--12. Defaults to a dry run that verifies the remapped state dict loads
cleanly (strict, zero missing/unexpected keys) before ever writing anything;
the original file is always backed up first.}

\func{\_remap\_key(key: str) -> str}
Applies each regex rule from the module's rename table in order, returning the
first match's replacement, or the key unchanged if it is already new-style
(the rules are mutually exclusive, so order does not affect correctness).

\func{\_is\_legacy\_checkpoint(state\_dict: dict) -> bool}
Returns \texttt{True} if any key still contains an old-style substring
exclusive to the legacy layout.

\func{\_migrate\_one(model\_name, run\_dir, pretrained, apply: bool) -> None}
Loads the checkpoint, skips (with a log message) if it is missing or already
new-style, otherwise remaps every key, then \textbf{verifies} the remap by
building a fresh model and diffing key sets --- raising rather than trusting an
unverified remap --- before ever writing. Only writes (backing up the original
first, and only if no backup already exists) when \texttt{apply=True}.

\func{main()}
Resolves which model names to migrate (default: the three known-affected
checkpoints), picks the correct pretrained-weights key per model, and calls
\texttt{\_migrate\_one} for each.

\begin{itemize}[nosep]
\cliarg{--paths}{config path}
\cliarg{--model-name}{repeatable; defaults to the 3 known-affected models}
\cliarg{--apply}{flag; without it, dry-run only}
\end{itemize}

\section{scripts/27\_evaluate\_test\_640\_models\_01\_to\_11.py}
\role{Computes test-set Dice (and extended metrics) for the five original
models always at a fixed 640x640 resolution --- ground truth is resized down by
the data loader, and predictions are \emph{never} upscaled back to native
camera resolution for comparison. This avoids the internal upscale step that
Ultralytics' own prediction path performs, giving a leaner ``what is the Dice
at 640, no upscaling'' number than the fuller Track A/B scripts.}

\func{\_load\_threshold(run\_dir) -> float} / \func{\_load\_threshold\_temp(run\_dir) -> tuple[float, float]}
The same validation-selected threshold (and, for YOLO, temperature) loading
pattern used throughout the project, raising rather than defaulting silently.

\func{\_score\_rows(rows, out\_dir, run\_name, best\_thr) -> None}
Writes the per-image CSV, computes the extended summary (2000 bootstrap
resamples), attaches run identity and an explicit note documenting the
no-upscale evaluation protocol, and writes the summary CSV.

\func{\_evaluate\_segformer(model\_name, pretrained\_key, paths, cfg, device, test\_df, out\_root) -> None}
Evaluates one SegFormer model over the test set at batch size 4 under autocast,
applying sigmoid and the validation threshold, and scoring with extended
per-image metrics.

\func{\_evaluate\_yolo(model\_name, paths, cfg, device, test\_df, out\_root) -> None}
Evaluates one YOLO model using \texttt{yolo\_raw\_class\_logits} resized to
match the batch's own 640 tensor size (never native resolution) plus
temperature-scaled probabilities --- bypassing Ultralytics' \texttt{.predict()}
entirely.

\func{main()}
Loops \texttt{\_evaluate\_segformer} over the 3 SegFormer models and
\texttt{\_evaluate\_yolo} over the 2 YOLO models.

\begin{itemize}[nosep]
\cliarg{--paths, --common}{config paths}
\end{itemize}
Module constants: \texttt{SURF\_DICE\_DELTA = 2.0}, \texttt{N\_BOOTSTRAP = 2000}, \texttt{EVAL\_BATCH = 4}.

\section{scripts/29\_benchmark\_inference\_all\_models.py}
\role{A GPU-only inference speed benchmark across every white-light-deployable
model, reusing the deployment-style timing logic later ported (and corrected)
into \texttt{pipeline/benchmark\_640.py}. Reports \texttt{raw\_forward\_fps}
(bare model call, isolating architecture speed) and \texttt{full\_pipeline\_fps}
(camera image in RAM through to a mask resized back to the \emph{original}
camera resolution --- the real deployment speed). Excludes disk read, model
loading, and warmup from timing. Note: this script's YOLO preprocessing scales
input to a plain 0--1 range with no ImageNet normalization, which
\texttt{pipeline/benchmark\_640.py}'s docstring identifies as inconsistent with
how YOLO's threshold and temperature were actually calibrated.}

\func{load\_image\_raw(image\_path: str) -> tuple[np.ndarray, tuple[int, int], str]}
Reads and BGR-to-RGB-converts an image, returning it with its native
\texttt{(height, width)} and filename stem. ``Disk read is setup, not part of
timed camera pipeline.''

\func{make\_segformer\_input(img\_rgb, size, device) -> torch.Tensor}
Resize plus ImageNet normalize plus to-tensor --- ``same inference
preprocessing structure as BruiseDataset for SegFormer.''

\func{make\_yolo\_input(img\_rgb, size, device) -> torch.Tensor}
Resize plus to-tensor, then manually divided by 255 --- \emph{not}
ImageNet-normalized, the inconsistency noted above.

\func{benchmark\_gpu\_fn(fn, device, warmup, n\_iters) -> tuple[float, float, float]}
Generic GPU timing harness: warms up untimed, times \texttt{n\_iters} calls
individually with synchronization bracketing each, returns mean/std/FPS.

\func{dtype\_configs(include\_fp32, include\_fp16)}
Returns the list of \texttt{(name, dtype, use\_amp)} tuples to benchmark,
depending on which precisions were requested.

\func{benchmark\_segformer\_one\_dtype(...) -> tuple[dict, list[dict]]}
Times a bare \texttt{raw\_forward} once on a warm input, then for
\emph{every} benchmark image times a \texttt{full\_pipeline} closure that
re-preprocesses, forwards, thresholds, and resizes the binary mask back to that
image's original resolution --- this per-image resize-back is the defining
feature of this script's ``full pipeline'' measurement (contrast with
\texttt{pipeline/benchmark\_640.py}, which keeps the mask at 640x640 by
design).

\func{benchmark\_yolo\_one\_dtype(...) -> tuple[dict, list[dict]]}
Same structure for YOLO, using \texttt{yolo\_raw\_class\_logits} plus
temperature-scaled probabilities before thresholding and resizing back to
native resolution.

\func{select\_benchmark\_images(paths, args) -> list[dict]}
Either benchmarks a single \texttt{--image}, or loads the fixed test manifest
(optionally shuffled by a fixed seed) and takes the first \texttt{--n-images}
rows, reading each one fully into memory as a record.

\func{main()}
Raises \texttt{RuntimeError} immediately if no CUDA device is available ---
``CPU FPS is not meaningful here.'' Loads/validates configs, selects benchmark
images, then loops the SegFormer and YOLO model lists (each optionally filtered
by \texttt{--models}), loading each (skip-with-warning on failure) and running
every resolution/dtype combination, writing the summary and per-image CSVs.

\begin{itemize}[nosep]
\cliarg{--paths, --common}{config paths}
\cliarg{--out, --per-image-out}{output CSV paths}
\cliarg{--image, --n-images, --sample-seed}{image selection (default 20 images, seed 2026)}
\cliarg{--resolutions}{default \texttt{[640]}}
\cliarg{--n-iters, --warmup}{default 200/20}
\cliarg{--fp32, --fp16}{flags; both run if neither given}
\cliarg{--models}{optional subset filter}
\end{itemize}

\section{scripts/31\_benchmark\_640\_local.py}
\role{A thin CLI wrapper around \texttt{pipeline.benchmark\_640.run\_benchmark\_640()}
--- the controlled 640x640 benchmark that runs on CPU or CUDA, unlike script 29
which hard-requires CUDA. \texttt{--device} defaults to auto-detection (CUDA if
available, else CPU) rather than raising, so it actually runs on this laptop.
Nearly all logic is deferred to \texttt{pipeline/benchmark\_640.py}; this file
only parses arguments and forwards them.}

\func{default\_device() -> str}
Returns \texttt{"cuda:0"} if available, else \texttt{"cpu"}.

\func{parse\_args() -> argparse.Namespace}
Defines the CLI surface (below).

\func{main()}
Forwards every parsed argument directly into \texttt{run\_benchmark\_640}.

\begin{itemize}[nosep]
\cliarg{--paths, --common, --models}{config paths (default \texttt{configs/paths.yaml}, \texttt{configs/common\_train.yaml}, \texttt{configs/benchmark\_640\_models.yaml})}
\cliarg{--out-dir}{default \texttt{results/benchmark\_640\_local}}
\cliarg{--device}{default: auto-detected}
\cliarg{--precision}{choices \texttt{fp32}/\texttt{fp16}, default \texttt{fp32}; fp16 auto-skipped off CUDA}
\cliarg{--max-images}{default 185 (the full fixed test set)}
\cliarg{--warmup, --repeats}{default 30/3}
\end{itemize}

\section{scripts/fix\_local\_manifest\_paths.py}
\role{A one-off utility rewriting hardcoded ORC-cluster Linux absolute paths
baked into manifest CSVs' \texttt{image\_path}/\texttt{mask\_path} columns, so
they resolve on this laptop. Necessary because \texttt{configs/paths.yaml} only
controls where the manifest \emph{files} are found --- it is never consulted
again once a manifest is loaded, since \texttt{BruiseDataset} reads whatever
string sits in those columns and hands it straight to \texttt{cv2.imread()}.
Only changes where a file is found on disk, never a pixel value or threshold,
so it cannot affect any already-recorded evaluation number.}

\func{remap\_one\_manifest(csv\_path: Path) -> None}
Backs up the untouched original to a sibling \texttt{.orc\_backup.csv} file (only
on first run), then always rebuilds the live CSV fresh from that backup with
every known ORC prefix replaced --- making the script idempotent and safe to
re-run any number of times, since the live file is always a correct rebuild
from the one true original, never a second pass over an already-replaced file.

\func{verify\_paths\_resolve(csv\_path: Path) -> list[str]}
Loads the CSV through \texttt{pipeline.data.normalize\_manifest} (reusing the
real pipeline's own column-normalization logic rather than re-guessing column
names) and returns every \texttt{image\_path}/\texttt{mask\_path} value that
does not resolve to a real file.

\func{main()}
For each of the 4 manifests wired into \texttt{paths.yaml} (the other manifests
in the training/test split and the two ALS manifests, plus the fixed test
manifest), skips with a warning if the file does not exist yet, otherwise
remaps and verifies it. Raises \texttt{RuntimeError} listing every path that
still fails to resolve after remapping.

\begin{itemize}[nosep]
\cliarg{--paths}{config path}
\end{itemize}

\section{scripts/visualize\_segformer\_stages.py}
\role{Produces a figure visualizing how a trained SegFormer-B2 checkpoint
processes one bruise image, stage by stage: input, patch embedding, each of the
four encoder stages, the decoded probability heatmap, and the final thresholded
mask. The docstring stresses these are the model's actual, real intermediate
activations on the given image (captured via a forward hook), not a schematic
illustration.}

\func{to\_heatmap(feature\_map: torch.Tensor, out\_h: int, out\_w: int) -> np.ndarray}
Averages a feature map across channels, min-max normalizes to $[0,1]$, and
resizes for smooth visualization.

\func{overlay(rgb\_uint8: np.ndarray, heat01: np.ndarray, alpha: float = 0.55) -> np.ndarray}
Converts a normalized heatmap to a colormap image and alpha-blends it over the
original photo.

\func{main()}
Builds the model and loads the checkpoint; registers a forward hook on the
first patch-embedding module to capture its raw spatial output; preprocesses
the input image by hand (resize, ImageNet normalize); runs one forward pass
capturing all four encoder stages' hidden states plus the final upsampled
logits; assembles a 7-panel figure (input, patch-embedding overlay, four stage
overlays with plain-English captions, the probability heatmap, and the binary
mask) and saves it as a PNG.

\begin{itemize}[nosep]
\cliarg{--checkpoint, --pretrained, --image}{required}
\cliarg{--img-size}{default 640}
\cliarg{--threshold}{default 0.5}
\cliarg{--out}{default \texttt{stage\_visualization.png}}
\end{itemize}

\section{scripts/yolo\_wl\_audit\_v1.py}
\role{A single-file, deliberately non-modular diagnostic script answering four
specific items raised in a stakeholder meeting, for the two white-light YOLO
models: (1) FPS at 640 only, with no upsample to native resolution --- reversing
earlier guidance that required timing the resize-back step; (2) test Dice at
640 with ground truth resized by the data loader, never native resolution;
(3) a temperature-scaling sanity check, reporting test Dice at both the
swept-best (temperature, threshold) and at temperature 1.0 side by side, so it
is visible whether temperature scaling actually helps; and (4) a comparison of
the project's own bare-PyTorch-plus-threshold pipeline against Ultralytics'
native prediction path on the same test set, at their respective native
resolutions (a difference the audit deliberately reports rather than silently
normalizing away).}

\func{\_make\_yolo\_input\_640(img\_rgb, img\_h, img\_w, device) -> torch.Tensor}
The same structural preprocessing as \texttt{BruiseDataset} for YOLO, kept
local rather than imported, per this project's convention that standalone
utility scripts do not import each other.

\func{\_benchmark\_fps\_640(yolo\_nn, img\_rgb, img\_h, img\_w, best\_thr, best\_temp, device) -> dict}
Times a raw-forward closure and a full-pipeline closure, both explicitly
without any resize back to camera resolution, per item 1 above.

\func{\_evaluate\_ultralytics\_native(wrapper, test\_df, img\_h, device\_str) -> pd.DataFrame}
Evaluates every test image using Ultralytics' own \texttt{.predict()} (its
internal argmax, resized to native ground-truth resolution) --- deliberately no
temperature scaling or threshold control, the ``as is'' side of item 4.

\func{\_audit\_one\_model(model\_name, paths, cfg, device) -> None}
Runs items 3, 2, 4, and 1 in sequence for one model: searches (temperature,
threshold) on validation and isolates the temperature-1.0 baseline; evaluates
the test set at 640 under both the swept-best and temperature-1.0 operating
points, logging a warning if temperature scaling does not actually help;
evaluates Ultralytics' native path at native resolution; benchmarks FPS at 640;
and writes one consolidated summary row combining every result.

\func{main()}
Hard-requires CUDA (``FPS numbers are meaningless on CPU'') and loops
\texttt{\_audit\_one\_model} over the direct and distilled YOLO models.

\begin{itemize}[nosep]
\cliarg{--paths, --common}{config paths}
\end{itemize}
Module constants: \texttt{N\_WARMUP = 20}, \texttt{N\_ITERS = 200}.

\end{document}
